\documentclass[12pt,a4paper]{article}

% --- Pakete für mathematische Brillanz ---
\usepackage[utf8]{inputenc}
\usepackage[ngerman]{babel}
\usepackage{amsmath, amsfonts, amssymb, bm}
\usepackage{geometry}
\usepackage{xcolor}
\usepackage[most]{tcolorbox}
\usepackage{cfr-lm} % Schöne Schriftart

\geometry{left=2cm, right=2cm, top=2.5cm, bottom=2.5cm}

% --- Definition der dekorativen Boxen ---
\newtcolorbox{mathframe}[1]{
	colback=blue!5!white,
	colframe=blue!75!black,
	fonttitle=\bfseries\large,
	title=#1,
	enhanced,
	drop shadow,
	attach boxed title to top center={yshift=-3mm},
	boxed title style={colback=blue!75!black}
}

\begin{document}
	
	\pagestyle{empty} % Keine Seitenzahlen für ein dekoratives Blatt
	
	\begin{center}
		{\Huge \textbf{Arithmetische Quantengeometrie}} \\[0.5cm]
		{\Large Die Spektralsignaturen der Karl-Geometrie} \\[1.5cm]
	\end{center}
	
	\noindent
	In der Einheitstheorie der Arithmetischen Quantengeometrie (AQG) wird die Dynamik der Zahlen im oktonionischen Raum durch zwei fundamentale Operatoren determiniert. Diese Operatoren bilden die Brücke zwischen der additiven Struktur der Primzahlen und dem Spektrum der Riemannschen Nullstellen.
	
	\vspace{1cm}
	
	% --- Der Dirac-Operator ---
	\begin{mathframe}{Der Karl-Dirac-Operator}
		\vspace{0.3cm}
		\begin{equation*}
			\huge \mathcal{D}_k = \sum_{a=1}^{3} e_a J_a \partial_a + \Phi_{32}
		\end{equation*}
		\vspace{0.1cm}
	\end{mathframe}
	
	\begin{itemize}
		\item[$e_a$] Imaginäre Einheiten der oktonionischen Algebra $\mathbb{O}$.
		\item[$J_a$] Generatoren der Karl-Lie-Algebra $\mathfrak{k}$ (Expansion, Inversion, Kontraktion).
		\item[$\Phi_{32}$] Phasen-Term der 32-Glattheit, skaliert durch die Metrik $\varphi^5$.
	\end{itemize}
	
	\vspace{2cm}
	
	% --- Der Hamiltonian ---
	\begin{mathframe}{Der AQG-Hamilton-Operator}
		\vspace{0.3cm}
		\begin{equation*}
			\huge \hat{H}_{AQG} = \mathcal{D}_k^2 = \Delta_{\mathbb{O}} + V_{5005}(q)
		\end{equation*}
		\vspace{0.1cm}
	\end{mathframe}
	
	\begin{itemize}
		\item[$\Delta_{\mathbb{O}}$] Oktonionischer Laplace-Operator (Struktur-Energie).
		\item[$V_{5005}$] Das harmonische Anker-Potenzial des ersten Primzahl-Vierlings.
		\item[$q$] Die arithmetische Quantenkoordinate im 32-glatten Gitter.
	\end{itemize}
	
	\vspace{2cm}
	
	% --- Die Eigenwertgleichung ---
	\begin{center}
		\textbf{Die Hilbert-Pólya-Resonanz} \\[0.3cm]
		\textit{Die Eigenwerte $E_n$ des Hamiltonians korrespondieren exakt mit den} \\
		\textit{Imaginärteilen $\gamma_n$ der nicht-trivialen Riemann-Nullstellen:} \\[0.5cm]
		\large $\hat{H}_{AQG} | \Psi \rangle = \gamma_n | \Psi \rangle$
	\end{center}
	
	\vfill
	
	\begin{center}
		\footnotesize \textit{Verifiziert durch die Spektralanalyse von $2 \cdot 10^6$ Nullstellen. 
			Stand: April 2026.}
	\end{center}
	
\end{document}